\section{Introduction}
\label{sec:introduction}
\textbf{Context and motivation.}
Introduce the research area in your own words. Describe a concrete problem,
who encounters it, and why solving it would matter. Define the scope so that
a reader outside your immediate specialty can follow the paper.

Describe a representative situation in which the problem becomes visible.
Explain what information is available, what decision or output is needed,
and what makes the task difficult.

\textbf{Gap and research question.}
Explain what existing approaches do not yet address. Support this gap with
real sources and formulate a specific, answerable research question.
Avoid claims of novelty that have not been checked against prior work.

Distinguish the practical challenge from the scientific question you intend
to investigate. Clarify whether the gap concerns a missing capability, an
untested assumption, or an insufficient explanation of existing behavior.
State the conditions under which answering your question would be useful,
and identify aspects of the broader problem that fall outside this paper.

\textbf{Key idea.}
Summarize your proposed approach in a few sentences. Explain why the idea
could address the gap and identify the assumption on which it depends.
Figure~\ref{fig:hook} highlights the central idea or motivating example.
Figure~\ref{fig:overview} provides a full-width overview of the method.

Explain the intuition connecting the proposed mechanism to the identified
gap. Contrast this intuition with a reasonable alternative and describe the
evidence needed to distinguish the two explanations. Introduce only the
technical terms necessary to understand the idea here; reserve detailed
notation and implementation choices for the later sections.

\textbf{Evaluation scope.}
Outline how the paper will examine its central claim. Identify the types
of comparisons, examples, and controlled changes needed for that purpose.
Explain how the evaluation relates to the motivating situation and what
it cannot establish. If the work is a proposal, clearly distinguish the
planned evaluation from any analysis already completed.

\textbf{Contributions.}
Replace the following items with concrete contributions supported by the paper:
\begin{itemize}
  \item A clearly scoped problem or research question: [replace].
  \item An approach or analysis addressing that question: [replace].
  \item Evidence and limitations from an appropriate evaluation: [replace].
\end{itemize}

The remainder of this paper reviews related work, defines the problem,
describes the method and evaluation, and discusses limitations.
